% Paste into an existing document, or use % Paste into an existing document, or use % Paste into an existing document, or use % Paste into an existing document, or use \input{ionisation}.
% Requires amsmath, amssymb, natbib, and references.bib.
\section{Ionisation processes in the solar atmosphere}
\label{sec:ionisation}

The first ionisation potential (FIP) effect depends on the different responses
of ions and neutral atoms to the ponderomotive force. We therefore require the
ionisation fraction of each element throughout the fractionation region.
This fraction depends on the local plasma conditions and on radiation received
from other parts of the atmosphere \citep{Laming2015}. The formulation below
describes a rate-based extension of the ionisation treatment in FIPpy.
It specifies the physical processes and approximations required for this
calculation; the existing Saha--CHIANTI interpolation is a separate prescription.

\subsection{Populations, rate coefficients, and transition rates}

For an element $X$ with atomic number $Z$, we define the charge-state fractions as
\begin{equation}
 f_{X,q} = \frac{n_{X^{q+}}}{n_X},
 \qquad n_X = \sum_{q=0}^{Z} n_{X^{q+}},
 \qquad \sum_{q=0}^{Z} f_{X,q}=1.
 \label{eq:ion-fractions}
\end{equation}
Here $q=0$ denotes the neutral stage. Each charge-state population includes
all of its bound levels. The total ionised fraction is
\begin{equation}
 \xi_X = \sum_{q=1}^{Z} f_{X,q} = 1-f_{X,0}.
 \label{eq:ion-total-fraction}
\end{equation}
When only neutral and singly ionised states are retained, this reduces to
$\xi_X=n_{X^+}/(n_{X^0}+n_{X^+})$. If higher stages are populated, their
contributions must also be included in the ionised fraction.

We distinguish a rate coefficient from a transition rate per particle.
An electron-impact coefficient $C_X$ has units of
$\mathrm{cm^3\,s^{-1}}$, while the corresponding transition rate is $n_e C_X$,
with units of $\mathrm{s^{-1}}$. The photoionisation rate $\Gamma_X$ already
has units of $\mathrm{s^{-1}}$. Multiplication by $n_e$ is required for
electron-collision coefficients and is absent from the photoionisation term.

\subsection{Electron-impact ionisation}

In direct ionisation, an incident free electron transfers enough energy to
remove a bound electron:
\begin{equation}
 X^{q+}+e^-\longrightarrow X^{(q+1)+}+2e^-.
 \label{eq:ion-direct}
\end{equation}
The coefficient depends on the electron energy distribution and on the
ionisation cross-section. For a thermal electron population, the coefficients
used here assume a Maxwellian distribution at electron temperature $T_e$.
An increase in $T_e$ changes the number of electrons above the threshold.

Excitation--autoionisation provides another electron-impact channel. An
electron collision excites the ion into an unstable state above an ionisation
threshold, after which the ion ejects an electron:
\begin{align}
 X^{q+}+e^- &\longrightarrow (X^{q+})^{**}+e^-,
 \\
 (X^{q+})^{**} &\longrightarrow X^{(q+1)+}+e^-.
\end{align}
The double asterisk denotes the unstable intermediate state. Ordinary
excitation to a bound level leaves the charge state unchanged; the
autoionisation step produces the increase in charge.

The electron-impact ionisation coefficient is therefore written as
\begin{equation}
 C_q(T_e)=C_{\mathrm{DI},q}(T_e)+C_{\mathrm{EA},q}(T_e),
 \label{eq:ion-electron-total}
\end{equation}
where DI and EA denote direct ionisation and excitation--autoionisation.
The rate per ion in stage $q$ is $n_e C_q$.

\subsection{Radiative and dielectronic recombination}

Radiative recombination occurs when an ion captures a free electron and emits
a photon:
\begin{equation}
 X^{(q+1)+}+e^-\longrightarrow X^{q+}+h\nu.
 \label{eq:ion-rr}
\end{equation}
The emitted photon carries away the energy released in the capture.
For the spontaneous radiative-recombination coefficients considered here,
the required inputs are the ion, the electron temperature, and the electron
density used to convert the coefficient into a rate. An imposed incident
radiation field is not required for this capture process.

Dielectronic recombination involves electron capture accompanied by excitation
of an electron already bound to the ion:
\begin{align}
 X^{(q+1)+}+e^- &\longrightarrow (X^{q+})^{**},
 \\
 (X^{q+})^{**} &\longrightarrow X^{q+}+h\nu.
 \label{eq:ion-dr}
\end{align}
The intermediate state can also autoionise. Recombination succeeds when
radiative stabilisation occurs before the captured electron is lost.
This mechanism requires an electron already bound to the initial ion.

In a model that represents both contributions through separate coefficients,
the total is
\begin{equation}
 \alpha_q(T_e)
 =\alpha_{\mathrm{RR},q}(T_e)+\alpha_{\mathrm{DR},q}(T_e),
 \label{eq:ion-recombination}
\end{equation}
where $\alpha_q$ describes recombination from stage $q$ to $q-1$.
The rate per ion is $n_e\alpha_q$. A database may instead supply a preferred
total fit that differs from the sum of its separately tabulated components.
The selected total and its provenance should then be retained consistently.

At finite electron density, collisions can reionise highly excited states
before stabilisation is complete. This suppresses the effective dielectronic
recombination coefficient. Density corrections or effective rates are required
when this process matters \citep{Nikolic2018,Dufresne2024}.
Applying a suppression factor to a coefficient that already includes this
effect would count the correction twice.

\subsection{Photoionisation}

A photon can remove a bound electron if its energy exceeds the relevant
ionisation threshold:
\begin{equation}
 X^{q+}+h\nu\longrightarrow X^{(q+1)+}+e^-.
 \label{eq:ion-photo-reaction}
\end{equation}
For a ground-state threshold $\chi_{X,q}$, the limiting frequency is
$\nu_{\mathrm{th},q}=\chi_{X,q}/h$. The rate per target particle is
\begin{equation}
 \Gamma_{X,q}(s)=4\pi\int_{\nu_{\mathrm{th},q}}^\infty
 \frac{J_\nu(s)}{h\nu}\,\sigma_{X,q}(\nu)\,\mathrm{d}\nu,
 \label{eq:ion-photo-frequency}
\end{equation}
where $J_\nu$ is the local mean intensity and $\sigma_{X,q}$ is the
photoionisation cross-section. The factor $J_\nu/(h\nu)$ converts energy
intensity into photon-number intensity. Integration over frequency and solid
angle gives a rate in $\mathrm{s^{-1}}$.

For a spectrum tabulated per unit wavelength, the equivalent expression is
\begin{equation}
 \Gamma_{X,q}(s)=4\pi\int_0^{\lambda_{\mathrm{th},q}}
 J_\lambda(s)\frac{\lambda}{hc}\,\sigma_{X,q}(\lambda)\,\mathrm{d}\lambda,
 \qquad \lambda_{\mathrm{th},q}=\frac{hc}{\chi_{X,q}}.
 \label{eq:ion-photo-wavelength}
\end{equation}
The spectral density and integration variable must use consistent units.
For example, $J_\lambda$ expressed per nanometre is integrated over
wavelength in nanometres, with the photon energy converted consistently.
Cross-sections vanish on the long-wavelength side of the threshold.
The integration grid must resolve the threshold and the spectral features
that contribute appreciably to the rate.

The analytic fits of \citet{Verner1996} provide a starting point for
ground-state cross-sections within their stated coverage. In a multilevel
model, excited levels have their own populations, thresholds, and
cross-sections. A ground-state rate then requires a population assumption
before it can represent the entire charge state.

\subsection{The radiation field and the role of opacity}

Radiation can reach a chromospheric parcel from the photosphere, the
transition region and corona, or emission within the chromosphere itself.
Photons from hotter regions can ionise atoms in a cooler parcel whose
local thermal electrons rarely exceed the threshold. The dominant source
depends on the element, wavelength, and position.

The specific intensity $I_\nu(s,\Omega)$ describes radiation arriving from a
direction $\Omega$. Its angular average is
\begin{equation}
 J_\nu(s)=\frac{1}{4\pi}\int I_\nu(s,\Omega)\,\mathrm{d}\Omega.
 \label{eq:ion-mean-intensity}
\end{equation}
Irradiance and mean intensity are different angular moments of the radiation
field. Converting an irradiance spectrum into $J_\nu$ therefore requires an
assumption about the emitting geometry.

For a uniformly bright photospheric surface, immediately above the surface,
the upward hemisphere gives $J_\nu=I_{\nu,\mathrm{ph}}/2$ and the emergent
surface flux is $F_{\nu,\mathrm{surf}}=\pi I_{\nu,\mathrm{ph}}$. Under this
assumption, an irradiance measured at distance $d$ gives
\begin{equation}
 J_{\nu,\mathrm{ph},0}
 =\frac{F_{\nu,d}}{2\pi}\left(\frac{d}{R_\odot}\right)^2.
 \label{eq:ion-irradiance-conversion}
\end{equation}
Downward coronal illumination requires its own geometry and emission model.
A spectrum measured at Earth does not uniquely determine the illumination
incident on a particular loop.

Along a ray with path coordinate $\ell$, radiative transfer satisfies
\begin{equation}
 \frac{\mathrm{d}I_\nu}{\mathrm{d}\ell}=-\kappa_\nu I_\nu+j_\nu,
 \label{eq:ion-transfer}
\end{equation}
where $\kappa_\nu$ is the extinction coefficient in $\mathrm{cm^{-1}}$ and
$j_\nu$ is the emission coefficient. Scattering contributes to extinction
and to the radiation redistributed into other directions and frequencies.
In the pure-absorption limit with no emission along the ray,
\begin{equation}
 I_\nu(\ell)=I_{\nu,0}\exp[-\tau_\nu(\ell)],
 \qquad
 \tau_\nu(\ell)=\int_{\ell_0}^{\ell}\kappa_\nu(\ell')\,\mathrm{d}\ell'.
 \label{eq:ion-attenuation}
\end{equation}

For ground-state neutral-hydrogen continuum absorption alone,
\begin{equation}
 \kappa_\nu=n_{\mathrm{H^0}}\sigma_{\mathrm{H^0}}(\nu),
 \qquad \tau_\nu=\sigma_{\mathrm{H^0}}(\nu)
          \int n_{\mathrm{H^0}}\,\mathrm{d}\ell.
 \label{eq:ion-h-opacity}
\end{equation}
Here the use of the total neutral density assumes that the ground state
dominates the absorbing neutral-H population.
Opacity determines the attenuation of an existing radiation field.
Calculating $\Gamma_X$ also requires the incident radiation or internal
emission and the target atom's cross-section.

Ground-state H continuum absorption vanishes below its ionisation threshold
of approximately $13.6\,\mathrm{eV}$. Metal continua and line opacity can still
affect radiation at these lower photon energies. In particular, resonant
absorption and scattering of Ly$\alpha$ require a line-transfer treatment
or a separately specified approximation. They are absent from
Equation~\eqref{eq:ion-h-opacity}. The assumed radiation and opacity treatment
should therefore accompany every calculated photoionisation rate
\citep{Laming2015,Dufresne2025}.

\subsection{Charge exchange with hydrogen}

Charge exchange transfers a bound electron between colliding particles:
\begin{equation}
 X^0+\mathrm{H^+}\rightleftharpoons X^++\mathrm{H^0}.
 \label{eq:ion-cx}
\end{equation}
In the forward direction, hydrogen captures an electron and $X$ becomes
ionised. In the reverse direction, $X^+$ captures hydrogen's electron and
becomes neutral. The electron remains bound to one of the reaction products,
and the total electric charge is conserved.

The corresponding rates per target particle are
\begin{equation}
 I_{\mathrm{cx},X}=n_{\mathrm{H^+}}k_{\mathrm{cx,ion},X},
 \qquad R_{\mathrm{cx},X}=n_{\mathrm{H^0}}k_{\mathrm{cx,rec},X}.
 \label{eq:ion-cx-rates}
\end{equation}
The coefficients have units of $\mathrm{cm^3\,s^{-1}}$ and depend on the
reaction and collision conditions. The appropriate heavy-particle
temperature need not equal $T_e$ in a plasma with different species
temperatures.

Oxygen and hydrogen have closely spaced first-ionisation thresholds.
Their nearly resonant charge exchange can consequently couple their
ionisation fractions when it dominates the other channels.
The fractions need not be exactly equal. Charge transfer can also affect
other elements, including sulphur, and requires species-specific data
\citep{Dufresne2021,Laming2021}. Elastic momentum-transfer collisions
enter the fractionation calculation separately; their collision frequencies
do not replace the reaction rates in Equation~\eqref{eq:ion-cx-rates}.

\subsection{Statistical equilibrium between neutral and singly ionised states}

For a two-stage model, the ionisation and neutralisation rates per particle are
\begin{align}
 I_X &= \Gamma_X+n_e C_X+n_{\mathrm{H^+}}k_{\mathrm{cx,ion},X},
 \label{eq:ion-forward-total}\\
 R_X &= n_e\alpha_X+n_{\mathrm{H^0}}k_{\mathrm{cx,rec},X}.
 \label{eq:ion-reverse-total}
\end{align}
Here $C_X$ describes $X^0\rightarrow X^+$, while $\alpha_X$ describes
$X^+\rightarrow X^0$. Statistical equilibrium requires
\begin{equation}
 n_{X^0}I_X=n_{X^+}R_X,
\end{equation}
which gives
\begin{equation}
 \xi_X=\frac{I_X}{I_X+R_X}.
 \label{eq:ion-two-state-equilibrium}
\end{equation}
This expression includes photoionisation, electron collisions, and charge
exchange. The term \emph{photoionisation equilibrium} is often used for
calculations including radiation, although it does not by itself establish
which channel dominates. The separate rates in
Equations~\eqref{eq:ion-forward-total} and \eqref{eq:ion-reverse-total}
allow that comparison to be made explicitly.

In the collision-only, two-stage limit,
\begin{equation}
 \xi_X^{\mathrm{CIE}}=\frac{C_X}{C_X+\alpha_X}.
 \label{eq:ion-two-state-cie}
\end{equation}
The common factor $n_e$ cancels when both coefficients depend only on
temperature. A multistage coronal calculation instead solves the coupled
population equations and returns the neutral fraction $f_{X,0}$.
Its total ionised fraction remains $1-f_{X,0}$.

\subsection{The LTE limit}

Under local thermodynamic equilibrium (LTE), the neutral-to-singly-ionised
ratio follows the Saha equation,
\begin{equation}
 \frac{n_{X^+}}{n_{X^0}}
 =\frac{2}{n_e}
  \left(\frac{2\pi m_e k_{\mathrm B}T}{h^2}\right)^{3/2}
  \frac{U_{X^+}(T)}{U_{X^0}(T)}
  \exp\left(-\frac{\chi_X}{k_{\mathrm B}T}\right),
 \label{eq:ion-saha}
\end{equation}
where $U(T)$ denotes a partition function. Replacing the partition functions
with constant statistical weights is an additional approximation.
Its adequacy depends on the population of excited levels.

LTE is a thermodynamic assumption, while statistical equilibrium specifies
that the population rates balance. Statistical equilibrium can hold in
non-LTE conditions when the radiation field and populations depart from their
local thermodynamic values. Recovering LTE from a rate solver requires the
relevant inverse processes and consistent detailed balance. An optically thin
ionisation and spontaneous-recombination network alone does not establish
the Saha limit.

\subsection{Extension to multiple stages and time dependence}

Let $I_q$ be the transition rate from stage $q$ to $q+1$, and $R_q$ the rate
from stage $q$ to $q-1$. With neighbouring-stage transitions, the local
population equations are
\begin{equation}
 \frac{\mathrm{d}f_{X,q}}{\mathrm{d}t}
 = I_{q-1}f_{X,q-1}+R_{q+1}f_{X,q+1}
   -(I_q+R_q)f_{X,q}.
 \label{eq:ion-multistage}
\end{equation}
Terms outside the retained stages are omitted, with $R_0=0$ and $I_Z=0$.
Statistical equilibrium sets the left-hand side to zero and enforces
Equation~\eqref{eq:ion-fractions}. A prescribed-background calculation holds
the atmospheric $T_e$, $n_e$, hydrogen populations, and specified radiation
fixed while solving the trace-element populations. Solving for $n_e$ through
charge conservation is a further coupling.

If conditions change faster than the populations relax, the time-dependent
equations must be evolved. For constant two-stage rates, the relaxation time
is $(I_X+R_X)^{-1}$. More general networks have several relaxation modes.
Hydrogen can retain a history of heating and shocks, which also changes
the collision partners and opacity seen by the metals \citep{Carlsson2002}.
Time-dependent calculations require initial populations and a sufficiently
resolved atmospheric and radiation history. Advection introduces transport
and boundary conditions in addition to the local reactions.

\subsection{Atomic coefficients from CHIANTI through fiasco}

The \texttt{fiasco} interface provides access to CHIANTI electron-impact
ionisation and recombination coefficients \citep{FiascoDocumentation}.
For a neutral--singly-ionised calculation, the required pair is the
\texttt{ionization\_rate} of the neutral ion object and the
\texttt{recombination\_rate} of the singly ionised object.
For example, Fe~I ionisation describes Fe~I to Fe~II, while Fe~II recombination
describes Fe~II to Fe~I. Fe~II ionisation instead describes Fe~II to Fe~III.
The two properties of a single ion object therefore refer to different
neighbouring-stage pairs.

The properties \texttt{direct\_ionization\_rate} and
\texttt{excitation\_autoionization\_rate} expose the ionisation components.
The properties \texttt{radiative\_recombination\_rate} and
\texttt{dielectronic\_recombination\_rate} expose recombination components.
For production calculations, the preferred total
\texttt{recombination\_rate} should be used because some ions have dedicated
total fits in CHIANTI \texttt{trparams} files.

Dataset availability must be checked for each ion. In the inspected
CHIANTI~11.0.2 Fe~II dataset, the \texttt{easplups} file is absent.
With fiasco~0.8.2, requesting the separate excitation--autoionisation
coefficient raises \texttt{MissingDatasetException}. Its total-ionisation
property handles this missing component internally and returns the
available direct-ionisation contribution. This is the database/software
treatment of the missing component; it does not demonstrate that the
physical process is identically zero.

Missing components should be recorded explicitly when plotting or exporting
the rates. The ordinary temperature-dependent properties also require
separate additions for the local photoionisation field, charge transfer,
and density effects appropriate to the selected model
\citep{Dufresne2024,Dufresne2025}.
.
% Requires amsmath, amssymb, natbib, and references.bib.
\section{Ionisation processes in the solar atmosphere}
\label{sec:ionisation}

The first ionisation potential (FIP) effect depends on the different responses
of ions and neutral atoms to the ponderomotive force. We therefore require the
ionisation fraction of each element throughout the fractionation region.
This fraction depends on the local plasma conditions and on radiation received
from other parts of the atmosphere \citep{Laming2015}. The formulation below
describes a rate-based extension of the ionisation treatment in FIPpy.
It specifies the physical processes and approximations required for this
calculation; the existing Saha--CHIANTI interpolation is a separate prescription.

\subsection{Populations, rate coefficients, and transition rates}

For an element $X$ with atomic number $Z$, we define the charge-state fractions as
\begin{equation}
 f_{X,q} = \frac{n_{X^{q+}}}{n_X},
 \qquad n_X = \sum_{q=0}^{Z} n_{X^{q+}},
 \qquad \sum_{q=0}^{Z} f_{X,q}=1.
 \label{eq:ion-fractions}
\end{equation}
Here $q=0$ denotes the neutral stage. Each charge-state population includes
all of its bound levels. The total ionised fraction is
\begin{equation}
 \xi_X = \sum_{q=1}^{Z} f_{X,q} = 1-f_{X,0}.
 \label{eq:ion-total-fraction}
\end{equation}
When only neutral and singly ionised states are retained, this reduces to
$\xi_X=n_{X^+}/(n_{X^0}+n_{X^+})$. If higher stages are populated, their
contributions must also be included in the ionised fraction.

We distinguish a rate coefficient from a transition rate per particle.
An electron-impact coefficient $C_X$ has units of
$\mathrm{cm^3\,s^{-1}}$, while the corresponding transition rate is $n_e C_X$,
with units of $\mathrm{s^{-1}}$. The photoionisation rate $\Gamma_X$ already
has units of $\mathrm{s^{-1}}$. Multiplication by $n_e$ is required for
electron-collision coefficients and is absent from the photoionisation term.

\subsection{Electron-impact ionisation}

In direct ionisation, an incident free electron transfers enough energy to
remove a bound electron:
\begin{equation}
 X^{q+}+e^-\longrightarrow X^{(q+1)+}+2e^-.
 \label{eq:ion-direct}
\end{equation}
The coefficient depends on the electron energy distribution and on the
ionisation cross-section. For a thermal electron population, the coefficients
used here assume a Maxwellian distribution at electron temperature $T_e$.
An increase in $T_e$ changes the number of electrons above the threshold.

Excitation--autoionisation provides another electron-impact channel. An
electron collision excites the ion into an unstable state above an ionisation
threshold, after which the ion ejects an electron:
\begin{align}
 X^{q+}+e^- &\longrightarrow (X^{q+})^{**}+e^-,
 \\
 (X^{q+})^{**} &\longrightarrow X^{(q+1)+}+e^-.
\end{align}
The double asterisk denotes the unstable intermediate state. Ordinary
excitation to a bound level leaves the charge state unchanged; the
autoionisation step produces the increase in charge.

The electron-impact ionisation coefficient is therefore written as
\begin{equation}
 C_q(T_e)=C_{\mathrm{DI},q}(T_e)+C_{\mathrm{EA},q}(T_e),
 \label{eq:ion-electron-total}
\end{equation}
where DI and EA denote direct ionisation and excitation--autoionisation.
The rate per ion in stage $q$ is $n_e C_q$.

\subsection{Radiative and dielectronic recombination}

Radiative recombination occurs when an ion captures a free electron and emits
a photon:
\begin{equation}
 X^{(q+1)+}+e^-\longrightarrow X^{q+}+h\nu.
 \label{eq:ion-rr}
\end{equation}
The emitted photon carries away the energy released in the capture.
For the spontaneous radiative-recombination coefficients considered here,
the required inputs are the ion, the electron temperature, and the electron
density used to convert the coefficient into a rate. An imposed incident
radiation field is not required for this capture process.

Dielectronic recombination involves electron capture accompanied by excitation
of an electron already bound to the ion:
\begin{align}
 X^{(q+1)+}+e^- &\longrightarrow (X^{q+})^{**},
 \\
 (X^{q+})^{**} &\longrightarrow X^{q+}+h\nu.
 \label{eq:ion-dr}
\end{align}
The intermediate state can also autoionise. Recombination succeeds when
radiative stabilisation occurs before the captured electron is lost.
This mechanism requires an electron already bound to the initial ion.

In a model that represents both contributions through separate coefficients,
the total is
\begin{equation}
 \alpha_q(T_e)
 =\alpha_{\mathrm{RR},q}(T_e)+\alpha_{\mathrm{DR},q}(T_e),
 \label{eq:ion-recombination}
\end{equation}
where $\alpha_q$ describes recombination from stage $q$ to $q-1$.
The rate per ion is $n_e\alpha_q$. A database may instead supply a preferred
total fit that differs from the sum of its separately tabulated components.
The selected total and its provenance should then be retained consistently.

At finite electron density, collisions can reionise highly excited states
before stabilisation is complete. This suppresses the effective dielectronic
recombination coefficient. Density corrections or effective rates are required
when this process matters \citep{Nikolic2018,Dufresne2024}.
Applying a suppression factor to a coefficient that already includes this
effect would count the correction twice.

\subsection{Photoionisation}

A photon can remove a bound electron if its energy exceeds the relevant
ionisation threshold:
\begin{equation}
 X^{q+}+h\nu\longrightarrow X^{(q+1)+}+e^-.
 \label{eq:ion-photo-reaction}
\end{equation}
For a ground-state threshold $\chi_{X,q}$, the limiting frequency is
$\nu_{\mathrm{th},q}=\chi_{X,q}/h$. The rate per target particle is
\begin{equation}
 \Gamma_{X,q}(s)=4\pi\int_{\nu_{\mathrm{th},q}}^\infty
 \frac{J_\nu(s)}{h\nu}\,\sigma_{X,q}(\nu)\,\mathrm{d}\nu,
 \label{eq:ion-photo-frequency}
\end{equation}
where $J_\nu$ is the local mean intensity and $\sigma_{X,q}$ is the
photoionisation cross-section. The factor $J_\nu/(h\nu)$ converts energy
intensity into photon-number intensity. Integration over frequency and solid
angle gives a rate in $\mathrm{s^{-1}}$.

For a spectrum tabulated per unit wavelength, the equivalent expression is
\begin{equation}
 \Gamma_{X,q}(s)=4\pi\int_0^{\lambda_{\mathrm{th},q}}
 J_\lambda(s)\frac{\lambda}{hc}\,\sigma_{X,q}(\lambda)\,\mathrm{d}\lambda,
 \qquad \lambda_{\mathrm{th},q}=\frac{hc}{\chi_{X,q}}.
 \label{eq:ion-photo-wavelength}
\end{equation}
The spectral density and integration variable must use consistent units.
For example, $J_\lambda$ expressed per nanometre is integrated over
wavelength in nanometres, with the photon energy converted consistently.
Cross-sections vanish on the long-wavelength side of the threshold.
The integration grid must resolve the threshold and the spectral features
that contribute appreciably to the rate.

The analytic fits of \citet{Verner1996} provide a starting point for
ground-state cross-sections within their stated coverage. In a multilevel
model, excited levels have their own populations, thresholds, and
cross-sections. A ground-state rate then requires a population assumption
before it can represent the entire charge state.

\subsection{The radiation field and the role of opacity}

Radiation can reach a chromospheric parcel from the photosphere, the
transition region and corona, or emission within the chromosphere itself.
Photons from hotter regions can ionise atoms in a cooler parcel whose
local thermal electrons rarely exceed the threshold. The dominant source
depends on the element, wavelength, and position.

The specific intensity $I_\nu(s,\Omega)$ describes radiation arriving from a
direction $\Omega$. Its angular average is
\begin{equation}
 J_\nu(s)=\frac{1}{4\pi}\int I_\nu(s,\Omega)\,\mathrm{d}\Omega.
 \label{eq:ion-mean-intensity}
\end{equation}
Irradiance and mean intensity are different angular moments of the radiation
field. Converting an irradiance spectrum into $J_\nu$ therefore requires an
assumption about the emitting geometry.

For a uniformly bright photospheric surface, immediately above the surface,
the upward hemisphere gives $J_\nu=I_{\nu,\mathrm{ph}}/2$ and the emergent
surface flux is $F_{\nu,\mathrm{surf}}=\pi I_{\nu,\mathrm{ph}}$. Under this
assumption, an irradiance measured at distance $d$ gives
\begin{equation}
 J_{\nu,\mathrm{ph},0}
 =\frac{F_{\nu,d}}{2\pi}\left(\frac{d}{R_\odot}\right)^2.
 \label{eq:ion-irradiance-conversion}
\end{equation}
Downward coronal illumination requires its own geometry and emission model.
A spectrum measured at Earth does not uniquely determine the illumination
incident on a particular loop.

Along a ray with path coordinate $\ell$, radiative transfer satisfies
\begin{equation}
 \frac{\mathrm{d}I_\nu}{\mathrm{d}\ell}=-\kappa_\nu I_\nu+j_\nu,
 \label{eq:ion-transfer}
\end{equation}
where $\kappa_\nu$ is the extinction coefficient in $\mathrm{cm^{-1}}$ and
$j_\nu$ is the emission coefficient. Scattering contributes to extinction
and to the radiation redistributed into other directions and frequencies.
In the pure-absorption limit with no emission along the ray,
\begin{equation}
 I_\nu(\ell)=I_{\nu,0}\exp[-\tau_\nu(\ell)],
 \qquad
 \tau_\nu(\ell)=\int_{\ell_0}^{\ell}\kappa_\nu(\ell')\,\mathrm{d}\ell'.
 \label{eq:ion-attenuation}
\end{equation}

For ground-state neutral-hydrogen continuum absorption alone,
\begin{equation}
 \kappa_\nu=n_{\mathrm{H^0}}\sigma_{\mathrm{H^0}}(\nu),
 \qquad \tau_\nu=\sigma_{\mathrm{H^0}}(\nu)
          \int n_{\mathrm{H^0}}\,\mathrm{d}\ell.
 \label{eq:ion-h-opacity}
\end{equation}
Here the use of the total neutral density assumes that the ground state
dominates the absorbing neutral-H population.
Opacity determines the attenuation of an existing radiation field.
Calculating $\Gamma_X$ also requires the incident radiation or internal
emission and the target atom's cross-section.

Ground-state H continuum absorption vanishes below its ionisation threshold
of approximately $13.6\,\mathrm{eV}$. Metal continua and line opacity can still
affect radiation at these lower photon energies. In particular, resonant
absorption and scattering of Ly$\alpha$ require a line-transfer treatment
or a separately specified approximation. They are absent from
Equation~\eqref{eq:ion-h-opacity}. The assumed radiation and opacity treatment
should therefore accompany every calculated photoionisation rate
\citep{Laming2015,Dufresne2025}.

\subsection{Charge exchange with hydrogen}

Charge exchange transfers a bound electron between colliding particles:
\begin{equation}
 X^0+\mathrm{H^+}\rightleftharpoons X^++\mathrm{H^0}.
 \label{eq:ion-cx}
\end{equation}
In the forward direction, hydrogen captures an electron and $X$ becomes
ionised. In the reverse direction, $X^+$ captures hydrogen's electron and
becomes neutral. The electron remains bound to one of the reaction products,
and the total electric charge is conserved.

The corresponding rates per target particle are
\begin{equation}
 I_{\mathrm{cx},X}=n_{\mathrm{H^+}}k_{\mathrm{cx,ion},X},
 \qquad R_{\mathrm{cx},X}=n_{\mathrm{H^0}}k_{\mathrm{cx,rec},X}.
 \label{eq:ion-cx-rates}
\end{equation}
The coefficients have units of $\mathrm{cm^3\,s^{-1}}$ and depend on the
reaction and collision conditions. The appropriate heavy-particle
temperature need not equal $T_e$ in a plasma with different species
temperatures.

Oxygen and hydrogen have closely spaced first-ionisation thresholds.
Their nearly resonant charge exchange can consequently couple their
ionisation fractions when it dominates the other channels.
The fractions need not be exactly equal. Charge transfer can also affect
other elements, including sulphur, and requires species-specific data
\citep{Dufresne2021,Laming2021}. Elastic momentum-transfer collisions
enter the fractionation calculation separately; their collision frequencies
do not replace the reaction rates in Equation~\eqref{eq:ion-cx-rates}.

\subsection{Statistical equilibrium between neutral and singly ionised states}

For a two-stage model, the ionisation and neutralisation rates per particle are
\begin{align}
 I_X &= \Gamma_X+n_e C_X+n_{\mathrm{H^+}}k_{\mathrm{cx,ion},X},
 \label{eq:ion-forward-total}\\
 R_X &= n_e\alpha_X+n_{\mathrm{H^0}}k_{\mathrm{cx,rec},X}.
 \label{eq:ion-reverse-total}
\end{align}
Here $C_X$ describes $X^0\rightarrow X^+$, while $\alpha_X$ describes
$X^+\rightarrow X^0$. Statistical equilibrium requires
\begin{equation}
 n_{X^0}I_X=n_{X^+}R_X,
\end{equation}
which gives
\begin{equation}
 \xi_X=\frac{I_X}{I_X+R_X}.
 \label{eq:ion-two-state-equilibrium}
\end{equation}
This expression includes photoionisation, electron collisions, and charge
exchange. The term \emph{photoionisation equilibrium} is often used for
calculations including radiation, although it does not by itself establish
which channel dominates. The separate rates in
Equations~\eqref{eq:ion-forward-total} and \eqref{eq:ion-reverse-total}
allow that comparison to be made explicitly.

In the collision-only, two-stage limit,
\begin{equation}
 \xi_X^{\mathrm{CIE}}=\frac{C_X}{C_X+\alpha_X}.
 \label{eq:ion-two-state-cie}
\end{equation}
The common factor $n_e$ cancels when both coefficients depend only on
temperature. A multistage coronal calculation instead solves the coupled
population equations and returns the neutral fraction $f_{X,0}$.
Its total ionised fraction remains $1-f_{X,0}$.

\subsection{The LTE limit}

Under local thermodynamic equilibrium (LTE), the neutral-to-singly-ionised
ratio follows the Saha equation,
\begin{equation}
 \frac{n_{X^+}}{n_{X^0}}
 =\frac{2}{n_e}
  \left(\frac{2\pi m_e k_{\mathrm B}T}{h^2}\right)^{3/2}
  \frac{U_{X^+}(T)}{U_{X^0}(T)}
  \exp\left(-\frac{\chi_X}{k_{\mathrm B}T}\right),
 \label{eq:ion-saha}
\end{equation}
where $U(T)$ denotes a partition function. Replacing the partition functions
with constant statistical weights is an additional approximation.
Its adequacy depends on the population of excited levels.

LTE is a thermodynamic assumption, while statistical equilibrium specifies
that the population rates balance. Statistical equilibrium can hold in
non-LTE conditions when the radiation field and populations depart from their
local thermodynamic values. Recovering LTE from a rate solver requires the
relevant inverse processes and consistent detailed balance. An optically thin
ionisation and spontaneous-recombination network alone does not establish
the Saha limit.

\subsection{Extension to multiple stages and time dependence}

Let $I_q$ be the transition rate from stage $q$ to $q+1$, and $R_q$ the rate
from stage $q$ to $q-1$. With neighbouring-stage transitions, the local
population equations are
\begin{equation}
 \frac{\mathrm{d}f_{X,q}}{\mathrm{d}t}
 = I_{q-1}f_{X,q-1}+R_{q+1}f_{X,q+1}
   -(I_q+R_q)f_{X,q}.
 \label{eq:ion-multistage}
\end{equation}
Terms outside the retained stages are omitted, with $R_0=0$ and $I_Z=0$.
Statistical equilibrium sets the left-hand side to zero and enforces
Equation~\eqref{eq:ion-fractions}. A prescribed-background calculation holds
the atmospheric $T_e$, $n_e$, hydrogen populations, and specified radiation
fixed while solving the trace-element populations. Solving for $n_e$ through
charge conservation is a further coupling.

If conditions change faster than the populations relax, the time-dependent
equations must be evolved. For constant two-stage rates, the relaxation time
is $(I_X+R_X)^{-1}$. More general networks have several relaxation modes.
Hydrogen can retain a history of heating and shocks, which also changes
the collision partners and opacity seen by the metals \citep{Carlsson2002}.
Time-dependent calculations require initial populations and a sufficiently
resolved atmospheric and radiation history. Advection introduces transport
and boundary conditions in addition to the local reactions.

\subsection{Atomic coefficients from CHIANTI through fiasco}

The \texttt{fiasco} interface provides access to CHIANTI electron-impact
ionisation and recombination coefficients \citep{FiascoDocumentation}.
For a neutral--singly-ionised calculation, the required pair is the
\texttt{ionization\_rate} of the neutral ion object and the
\texttt{recombination\_rate} of the singly ionised object.
For example, Fe~I ionisation describes Fe~I to Fe~II, while Fe~II recombination
describes Fe~II to Fe~I. Fe~II ionisation instead describes Fe~II to Fe~III.
The two properties of a single ion object therefore refer to different
neighbouring-stage pairs.

The properties \texttt{direct\_ionization\_rate} and
\texttt{excitation\_autoionization\_rate} expose the ionisation components.
The properties \texttt{radiative\_recombination\_rate} and
\texttt{dielectronic\_recombination\_rate} expose recombination components.
For production calculations, the preferred total
\texttt{recombination\_rate} should be used because some ions have dedicated
total fits in CHIANTI \texttt{trparams} files.

Dataset availability must be checked for each ion. In the inspected
CHIANTI~11.0.2 Fe~II dataset, the \texttt{easplups} file is absent.
With fiasco~0.8.2, requesting the separate excitation--autoionisation
coefficient raises \texttt{MissingDatasetException}. Its total-ionisation
property handles this missing component internally and returns the
available direct-ionisation contribution. This is the database/software
treatment of the missing component; it does not demonstrate that the
physical process is identically zero.

Missing components should be recorded explicitly when plotting or exporting
the rates. The ordinary temperature-dependent properties also require
separate additions for the local photoionisation field, charge transfer,
and density effects appropriate to the selected model
\citep{Dufresne2024,Dufresne2025}.
.
% Requires amsmath, amssymb, natbib, and references.bib.
\section{Ionisation processes in the solar atmosphere}
\label{sec:ionisation}

The first ionisation potential (FIP) effect depends on the different responses
of ions and neutral atoms to the ponderomotive force. We therefore require the
ionisation fraction of each element throughout the fractionation region.
This fraction depends on the local plasma conditions and on radiation received
from other parts of the atmosphere \citep{Laming2015}. The formulation below
describes a rate-based extension of the ionisation treatment in FIPpy.
It specifies the physical processes and approximations required for this
calculation; the existing Saha--CHIANTI interpolation is a separate prescription.

\subsection{Populations, rate coefficients, and transition rates}

For an element $X$ with atomic number $Z$, we define the charge-state fractions as
\begin{equation}
 f_{X,q} = \frac{n_{X^{q+}}}{n_X},
 \qquad n_X = \sum_{q=0}^{Z} n_{X^{q+}},
 \qquad \sum_{q=0}^{Z} f_{X,q}=1.
 \label{eq:ion-fractions}
\end{equation}
Here $q=0$ denotes the neutral stage. Each charge-state population includes
all of its bound levels. The total ionised fraction is
\begin{equation}
 \xi_X = \sum_{q=1}^{Z} f_{X,q} = 1-f_{X,0}.
 \label{eq:ion-total-fraction}
\end{equation}
When only neutral and singly ionised states are retained, this reduces to
$\xi_X=n_{X^+}/(n_{X^0}+n_{X^+})$. If higher stages are populated, their
contributions must also be included in the ionised fraction.

We distinguish a rate coefficient from a transition rate per particle.
An electron-impact coefficient $C_X$ has units of
$\mathrm{cm^3\,s^{-1}}$, while the corresponding transition rate is $n_e C_X$,
with units of $\mathrm{s^{-1}}$. The photoionisation rate $\Gamma_X$ already
has units of $\mathrm{s^{-1}}$. Multiplication by $n_e$ is required for
electron-collision coefficients and is absent from the photoionisation term.

\subsection{Electron-impact ionisation}

In direct ionisation, an incident free electron transfers enough energy to
remove a bound electron:
\begin{equation}
 X^{q+}+e^-\longrightarrow X^{(q+1)+}+2e^-.
 \label{eq:ion-direct}
\end{equation}
The coefficient depends on the electron energy distribution and on the
ionisation cross-section. For a thermal electron population, the coefficients
used here assume a Maxwellian distribution at electron temperature $T_e$.
An increase in $T_e$ changes the number of electrons above the threshold.

Excitation--autoionisation provides another electron-impact channel. An
electron collision excites the ion into an unstable state above an ionisation
threshold, after which the ion ejects an electron:
\begin{align}
 X^{q+}+e^- &\longrightarrow (X^{q+})^{**}+e^-,
 \\
 (X^{q+})^{**} &\longrightarrow X^{(q+1)+}+e^-.
\end{align}
The double asterisk denotes the unstable intermediate state. Ordinary
excitation to a bound level leaves the charge state unchanged; the
autoionisation step produces the increase in charge.

The electron-impact ionisation coefficient is therefore written as
\begin{equation}
 C_q(T_e)=C_{\mathrm{DI},q}(T_e)+C_{\mathrm{EA},q}(T_e),
 \label{eq:ion-electron-total}
\end{equation}
where DI and EA denote direct ionisation and excitation--autoionisation.
The rate per ion in stage $q$ is $n_e C_q$.

\subsection{Radiative and dielectronic recombination}

Radiative recombination occurs when an ion captures a free electron and emits
a photon:
\begin{equation}
 X^{(q+1)+}+e^-\longrightarrow X^{q+}+h\nu.
 \label{eq:ion-rr}
\end{equation}
The emitted photon carries away the energy released in the capture.
For the spontaneous radiative-recombination coefficients considered here,
the required inputs are the ion, the electron temperature, and the electron
density used to convert the coefficient into a rate. An imposed incident
radiation field is not required for this capture process.

Dielectronic recombination involves electron capture accompanied by excitation
of an electron already bound to the ion:
\begin{align}
 X^{(q+1)+}+e^- &\longrightarrow (X^{q+})^{**},
 \\
 (X^{q+})^{**} &\longrightarrow X^{q+}+h\nu.
 \label{eq:ion-dr}
\end{align}
The intermediate state can also autoionise. Recombination succeeds when
radiative stabilisation occurs before the captured electron is lost.
This mechanism requires an electron already bound to the initial ion.

In a model that represents both contributions through separate coefficients,
the total is
\begin{equation}
 \alpha_q(T_e)
 =\alpha_{\mathrm{RR},q}(T_e)+\alpha_{\mathrm{DR},q}(T_e),
 \label{eq:ion-recombination}
\end{equation}
where $\alpha_q$ describes recombination from stage $q$ to $q-1$.
The rate per ion is $n_e\alpha_q$. A database may instead supply a preferred
total fit that differs from the sum of its separately tabulated components.
The selected total and its provenance should then be retained consistently.

At finite electron density, collisions can reionise highly excited states
before stabilisation is complete. This suppresses the effective dielectronic
recombination coefficient. Density corrections or effective rates are required
when this process matters \citep{Nikolic2018,Dufresne2024}.
Applying a suppression factor to a coefficient that already includes this
effect would count the correction twice.

\subsection{Photoionisation}

A photon can remove a bound electron if its energy exceeds the relevant
ionisation threshold:
\begin{equation}
 X^{q+}+h\nu\longrightarrow X^{(q+1)+}+e^-.
 \label{eq:ion-photo-reaction}
\end{equation}
For a ground-state threshold $\chi_{X,q}$, the limiting frequency is
$\nu_{\mathrm{th},q}=\chi_{X,q}/h$. The rate per target particle is
\begin{equation}
 \Gamma_{X,q}(s)=4\pi\int_{\nu_{\mathrm{th},q}}^\infty
 \frac{J_\nu(s)}{h\nu}\,\sigma_{X,q}(\nu)\,\mathrm{d}\nu,
 \label{eq:ion-photo-frequency}
\end{equation}
where $J_\nu$ is the local mean intensity and $\sigma_{X,q}$ is the
photoionisation cross-section. The factor $J_\nu/(h\nu)$ converts energy
intensity into photon-number intensity. Integration over frequency and solid
angle gives a rate in $\mathrm{s^{-1}}$.

For a spectrum tabulated per unit wavelength, the equivalent expression is
\begin{equation}
 \Gamma_{X,q}(s)=4\pi\int_0^{\lambda_{\mathrm{th},q}}
 J_\lambda(s)\frac{\lambda}{hc}\,\sigma_{X,q}(\lambda)\,\mathrm{d}\lambda,
 \qquad \lambda_{\mathrm{th},q}=\frac{hc}{\chi_{X,q}}.
 \label{eq:ion-photo-wavelength}
\end{equation}
The spectral density and integration variable must use consistent units.
For example, $J_\lambda$ expressed per nanometre is integrated over
wavelength in nanometres, with the photon energy converted consistently.
Cross-sections vanish on the long-wavelength side of the threshold.
The integration grid must resolve the threshold and the spectral features
that contribute appreciably to the rate.

The analytic fits of \citet{Verner1996} provide a starting point for
ground-state cross-sections within their stated coverage. In a multilevel
model, excited levels have their own populations, thresholds, and
cross-sections. A ground-state rate then requires a population assumption
before it can represent the entire charge state.

\subsection{The radiation field and the role of opacity}

Radiation can reach a chromospheric parcel from the photosphere, the
transition region and corona, or emission within the chromosphere itself.
Photons from hotter regions can ionise atoms in a cooler parcel whose
local thermal electrons rarely exceed the threshold. The dominant source
depends on the element, wavelength, and position.

The specific intensity $I_\nu(s,\Omega)$ describes radiation arriving from a
direction $\Omega$. Its angular average is
\begin{equation}
 J_\nu(s)=\frac{1}{4\pi}\int I_\nu(s,\Omega)\,\mathrm{d}\Omega.
 \label{eq:ion-mean-intensity}
\end{equation}
Irradiance and mean intensity are different angular moments of the radiation
field. Converting an irradiance spectrum into $J_\nu$ therefore requires an
assumption about the emitting geometry.

For a uniformly bright photospheric surface, immediately above the surface,
the upward hemisphere gives $J_\nu=I_{\nu,\mathrm{ph}}/2$ and the emergent
surface flux is $F_{\nu,\mathrm{surf}}=\pi I_{\nu,\mathrm{ph}}$. Under this
assumption, an irradiance measured at distance $d$ gives
\begin{equation}
 J_{\nu,\mathrm{ph},0}
 =\frac{F_{\nu,d}}{2\pi}\left(\frac{d}{R_\odot}\right)^2.
 \label{eq:ion-irradiance-conversion}
\end{equation}
Downward coronal illumination requires its own geometry and emission model.
A spectrum measured at Earth does not uniquely determine the illumination
incident on a particular loop.

Along a ray with path coordinate $\ell$, radiative transfer satisfies
\begin{equation}
 \frac{\mathrm{d}I_\nu}{\mathrm{d}\ell}=-\kappa_\nu I_\nu+j_\nu,
 \label{eq:ion-transfer}
\end{equation}
where $\kappa_\nu$ is the extinction coefficient in $\mathrm{cm^{-1}}$ and
$j_\nu$ is the emission coefficient. Scattering contributes to extinction
and to the radiation redistributed into other directions and frequencies.
In the pure-absorption limit with no emission along the ray,
\begin{equation}
 I_\nu(\ell)=I_{\nu,0}\exp[-\tau_\nu(\ell)],
 \qquad
 \tau_\nu(\ell)=\int_{\ell_0}^{\ell}\kappa_\nu(\ell')\,\mathrm{d}\ell'.
 \label{eq:ion-attenuation}
\end{equation}

For ground-state neutral-hydrogen continuum absorption alone,
\begin{equation}
 \kappa_\nu=n_{\mathrm{H^0}}\sigma_{\mathrm{H^0}}(\nu),
 \qquad \tau_\nu=\sigma_{\mathrm{H^0}}(\nu)
          \int n_{\mathrm{H^0}}\,\mathrm{d}\ell.
 \label{eq:ion-h-opacity}
\end{equation}
Here the use of the total neutral density assumes that the ground state
dominates the absorbing neutral-H population.
Opacity determines the attenuation of an existing radiation field.
Calculating $\Gamma_X$ also requires the incident radiation or internal
emission and the target atom's cross-section.

Ground-state H continuum absorption vanishes below its ionisation threshold
of approximately $13.6\,\mathrm{eV}$. Metal continua and line opacity can still
affect radiation at these lower photon energies. In particular, resonant
absorption and scattering of Ly$\alpha$ require a line-transfer treatment
or a separately specified approximation. They are absent from
Equation~\eqref{eq:ion-h-opacity}. The assumed radiation and opacity treatment
should therefore accompany every calculated photoionisation rate
\citep{Laming2015,Dufresne2025}.

\subsection{Charge exchange with hydrogen}

Charge exchange transfers a bound electron between colliding particles:
\begin{equation}
 X^0+\mathrm{H^+}\rightleftharpoons X^++\mathrm{H^0}.
 \label{eq:ion-cx}
\end{equation}
In the forward direction, hydrogen captures an electron and $X$ becomes
ionised. In the reverse direction, $X^+$ captures hydrogen's electron and
becomes neutral. The electron remains bound to one of the reaction products,
and the total electric charge is conserved.

The corresponding rates per target particle are
\begin{equation}
 I_{\mathrm{cx},X}=n_{\mathrm{H^+}}k_{\mathrm{cx,ion},X},
 \qquad R_{\mathrm{cx},X}=n_{\mathrm{H^0}}k_{\mathrm{cx,rec},X}.
 \label{eq:ion-cx-rates}
\end{equation}
The coefficients have units of $\mathrm{cm^3\,s^{-1}}$ and depend on the
reaction and collision conditions. The appropriate heavy-particle
temperature need not equal $T_e$ in a plasma with different species
temperatures.

Oxygen and hydrogen have closely spaced first-ionisation thresholds.
Their nearly resonant charge exchange can consequently couple their
ionisation fractions when it dominates the other channels.
The fractions need not be exactly equal. Charge transfer can also affect
other elements, including sulphur, and requires species-specific data
\citep{Dufresne2021,Laming2021}. Elastic momentum-transfer collisions
enter the fractionation calculation separately; their collision frequencies
do not replace the reaction rates in Equation~\eqref{eq:ion-cx-rates}.

\subsection{Statistical equilibrium between neutral and singly ionised states}

For a two-stage model, the ionisation and neutralisation rates per particle are
\begin{align}
 I_X &= \Gamma_X+n_e C_X+n_{\mathrm{H^+}}k_{\mathrm{cx,ion},X},
 \label{eq:ion-forward-total}\\
 R_X &= n_e\alpha_X+n_{\mathrm{H^0}}k_{\mathrm{cx,rec},X}.
 \label{eq:ion-reverse-total}
\end{align}
Here $C_X$ describes $X^0\rightarrow X^+$, while $\alpha_X$ describes
$X^+\rightarrow X^0$. Statistical equilibrium requires
\begin{equation}
 n_{X^0}I_X=n_{X^+}R_X,
\end{equation}
which gives
\begin{equation}
 \xi_X=\frac{I_X}{I_X+R_X}.
 \label{eq:ion-two-state-equilibrium}
\end{equation}
This expression includes photoionisation, electron collisions, and charge
exchange. The term \emph{photoionisation equilibrium} is often used for
calculations including radiation, although it does not by itself establish
which channel dominates. The separate rates in
Equations~\eqref{eq:ion-forward-total} and \eqref{eq:ion-reverse-total}
allow that comparison to be made explicitly.

In the collision-only, two-stage limit,
\begin{equation}
 \xi_X^{\mathrm{CIE}}=\frac{C_X}{C_X+\alpha_X}.
 \label{eq:ion-two-state-cie}
\end{equation}
The common factor $n_e$ cancels when both coefficients depend only on
temperature. A multistage coronal calculation instead solves the coupled
population equations and returns the neutral fraction $f_{X,0}$.
Its total ionised fraction remains $1-f_{X,0}$.

\subsection{The LTE limit}

Under local thermodynamic equilibrium (LTE), the neutral-to-singly-ionised
ratio follows the Saha equation,
\begin{equation}
 \frac{n_{X^+}}{n_{X^0}}
 =\frac{2}{n_e}
  \left(\frac{2\pi m_e k_{\mathrm B}T}{h^2}\right)^{3/2}
  \frac{U_{X^+}(T)}{U_{X^0}(T)}
  \exp\left(-\frac{\chi_X}{k_{\mathrm B}T}\right),
 \label{eq:ion-saha}
\end{equation}
where $U(T)$ denotes a partition function. Replacing the partition functions
with constant statistical weights is an additional approximation.
Its adequacy depends on the population of excited levels.

LTE is a thermodynamic assumption, while statistical equilibrium specifies
that the population rates balance. Statistical equilibrium can hold in
non-LTE conditions when the radiation field and populations depart from their
local thermodynamic values. Recovering LTE from a rate solver requires the
relevant inverse processes and consistent detailed balance. An optically thin
ionisation and spontaneous-recombination network alone does not establish
the Saha limit.

\subsection{Extension to multiple stages and time dependence}

Let $I_q$ be the transition rate from stage $q$ to $q+1$, and $R_q$ the rate
from stage $q$ to $q-1$. With neighbouring-stage transitions, the local
population equations are
\begin{equation}
 \frac{\mathrm{d}f_{X,q}}{\mathrm{d}t}
 = I_{q-1}f_{X,q-1}+R_{q+1}f_{X,q+1}
   -(I_q+R_q)f_{X,q}.
 \label{eq:ion-multistage}
\end{equation}
Terms outside the retained stages are omitted, with $R_0=0$ and $I_Z=0$.
Statistical equilibrium sets the left-hand side to zero and enforces
Equation~\eqref{eq:ion-fractions}. A prescribed-background calculation holds
the atmospheric $T_e$, $n_e$, hydrogen populations, and specified radiation
fixed while solving the trace-element populations. Solving for $n_e$ through
charge conservation is a further coupling.

If conditions change faster than the populations relax, the time-dependent
equations must be evolved. For constant two-stage rates, the relaxation time
is $(I_X+R_X)^{-1}$. More general networks have several relaxation modes.
Hydrogen can retain a history of heating and shocks, which also changes
the collision partners and opacity seen by the metals \citep{Carlsson2002}.
Time-dependent calculations require initial populations and a sufficiently
resolved atmospheric and radiation history. Advection introduces transport
and boundary conditions in addition to the local reactions.

\subsection{Atomic coefficients from CHIANTI through fiasco}

The \texttt{fiasco} interface provides access to CHIANTI electron-impact
ionisation and recombination coefficients \citep{FiascoDocumentation}.
For a neutral--singly-ionised calculation, the required pair is the
\texttt{ionization\_rate} of the neutral ion object and the
\texttt{recombination\_rate} of the singly ionised object.
For example, Fe~I ionisation describes Fe~I to Fe~II, while Fe~II recombination
describes Fe~II to Fe~I. Fe~II ionisation instead describes Fe~II to Fe~III.
The two properties of a single ion object therefore refer to different
neighbouring-stage pairs.

The properties \texttt{direct\_ionization\_rate} and
\texttt{excitation\_autoionization\_rate} expose the ionisation components.
The properties \texttt{radiative\_recombination\_rate} and
\texttt{dielectronic\_recombination\_rate} expose recombination components.
For production calculations, the preferred total
\texttt{recombination\_rate} should be used because some ions have dedicated
total fits in CHIANTI \texttt{trparams} files.

Dataset availability must be checked for each ion. In the inspected
CHIANTI~11.0.2 Fe~II dataset, the \texttt{easplups} file is absent.
With fiasco~0.8.2, requesting the separate excitation--autoionisation
coefficient raises \texttt{MissingDatasetException}. Its total-ionisation
property handles this missing component internally and returns the
available direct-ionisation contribution. This is the database/software
treatment of the missing component; it does not demonstrate that the
physical process is identically zero.

Missing components should be recorded explicitly when plotting or exporting
the rates. The ordinary temperature-dependent properties also require
separate additions for the local photoionisation field, charge transfer,
and density effects appropriate to the selected model
\citep{Dufresne2024,Dufresne2025}.
.
% Requires amsmath, amssymb, natbib, and references.bib.
\section{Ionisation processes in the solar atmosphere}
\label{sec:ionisation}

The first ionisation potential (FIP) effect depends on the different responses
of ions and neutral atoms to the ponderomotive force. We therefore require the
ionisation fraction of each element throughout the fractionation region.
This fraction depends on the local plasma conditions and on radiation received
from other parts of the atmosphere \citep{Laming2015}. The formulation below
describes a rate-based extension of the ionisation treatment in FIPpy.
It specifies the physical processes and approximations required for this
calculation; the existing Saha--CHIANTI interpolation is a separate prescription.

\subsection{Populations, rate coefficients, and transition rates}

For an element $X$ with atomic number $Z$, we define the charge-state fractions as
\begin{equation}
 f_{X,q} = \frac{n_{X^{q+}}}{n_X},
 \qquad n_X = \sum_{q=0}^{Z} n_{X^{q+}},
 \qquad \sum_{q=0}^{Z} f_{X,q}=1.
 \label{eq:ion-fractions}
\end{equation}
Here $q=0$ denotes the neutral stage. Each charge-state population includes
all of its bound levels. The total ionised fraction is
\begin{equation}
 \xi_X = \sum_{q=1}^{Z} f_{X,q} = 1-f_{X,0}.
 \label{eq:ion-total-fraction}
\end{equation}
When only neutral and singly ionised states are retained, this reduces to
$\xi_X=n_{X^+}/(n_{X^0}+n_{X^+})$. If higher stages are populated, their
contributions must also be included in the ionised fraction.

We distinguish a rate coefficient from a transition rate per particle.
An electron-impact coefficient $C_X$ has units of
$\mathrm{cm^3\,s^{-1}}$, while the corresponding transition rate is $n_e C_X$,
with units of $\mathrm{s^{-1}}$. The photoionisation rate $\Gamma_X$ already
has units of $\mathrm{s^{-1}}$. Multiplication by $n_e$ is required for
electron-collision coefficients and is absent from the photoionisation term.

\subsection{Electron-impact ionisation}

In direct ionisation, an incident free electron transfers enough energy to
remove a bound electron:
\begin{equation}
 X^{q+}+e^-\longrightarrow X^{(q+1)+}+2e^-.
 \label{eq:ion-direct}
\end{equation}
The coefficient depends on the electron energy distribution and on the
ionisation cross-section. For a thermal electron population, the coefficients
used here assume a Maxwellian distribution at electron temperature $T_e$.
An increase in $T_e$ changes the number of electrons above the threshold.

Excitation--autoionisation provides another electron-impact channel. An
electron collision excites the ion into an unstable state above an ionisation
threshold, after which the ion ejects an electron:
\begin{align}
 X^{q+}+e^- &\longrightarrow (X^{q+})^{**}+e^-,
 \\
 (X^{q+})^{**} &\longrightarrow X^{(q+1)+}+e^-.
\end{align}
The double asterisk denotes the unstable intermediate state. Ordinary
excitation to a bound level leaves the charge state unchanged; the
autoionisation step produces the increase in charge.

The electron-impact ionisation coefficient is therefore written as
\begin{equation}
 C_q(T_e)=C_{\mathrm{DI},q}(T_e)+C_{\mathrm{EA},q}(T_e),
 \label{eq:ion-electron-total}
\end{equation}
where DI and EA denote direct ionisation and excitation--autoionisation.
The rate per ion in stage $q$ is $n_e C_q$.

\subsection{Radiative and dielectronic recombination}

Radiative recombination occurs when an ion captures a free electron and emits
a photon:
\begin{equation}
 X^{(q+1)+}+e^-\longrightarrow X^{q+}+h\nu.
 \label{eq:ion-rr}
\end{equation}
The emitted photon carries away the energy released in the capture.
For the spontaneous radiative-recombination coefficients considered here,
the required inputs are the ion, the electron temperature, and the electron
density used to convert the coefficient into a rate. An imposed incident
radiation field is not required for this capture process.

Dielectronic recombination involves electron capture accompanied by excitation
of an electron already bound to the ion:
\begin{align}
 X^{(q+1)+}+e^- &\longrightarrow (X^{q+})^{**},
 \\
 (X^{q+})^{**} &\longrightarrow X^{q+}+h\nu.
 \label{eq:ion-dr}
\end{align}
The intermediate state can also autoionise. Recombination succeeds when
radiative stabilisation occurs before the captured electron is lost.
This mechanism requires an electron already bound to the initial ion.

In a model that represents both contributions through separate coefficients,
the total is
\begin{equation}
 \alpha_q(T_e)
 =\alpha_{\mathrm{RR},q}(T_e)+\alpha_{\mathrm{DR},q}(T_e),
 \label{eq:ion-recombination}
\end{equation}
where $\alpha_q$ describes recombination from stage $q$ to $q-1$.
The rate per ion is $n_e\alpha_q$. A database may instead supply a preferred
total fit that differs from the sum of its separately tabulated components.
The selected total and its provenance should then be retained consistently.

At finite electron density, collisions can reionise highly excited states
before stabilisation is complete. This suppresses the effective dielectronic
recombination coefficient. Density corrections or effective rates are required
when this process matters \citep{Nikolic2018,Dufresne2024}.
Applying a suppression factor to a coefficient that already includes this
effect would count the correction twice.

\subsection{Photoionisation}

A photon can remove a bound electron if its energy exceeds the relevant
ionisation threshold:
\begin{equation}
 X^{q+}+h\nu\longrightarrow X^{(q+1)+}+e^-.
 \label{eq:ion-photo-reaction}
\end{equation}
For a ground-state threshold $\chi_{X,q}$, the limiting frequency is
$\nu_{\mathrm{th},q}=\chi_{X,q}/h$. The rate per target particle is
\begin{equation}
 \Gamma_{X,q}(s)=4\pi\int_{\nu_{\mathrm{th},q}}^\infty
 \frac{J_\nu(s)}{h\nu}\,\sigma_{X,q}(\nu)\,\mathrm{d}\nu,
 \label{eq:ion-photo-frequency}
\end{equation}
where $J_\nu$ is the local mean intensity and $\sigma_{X,q}$ is the
photoionisation cross-section. The factor $J_\nu/(h\nu)$ converts energy
intensity into photon-number intensity. Integration over frequency and solid
angle gives a rate in $\mathrm{s^{-1}}$.

For a spectrum tabulated per unit wavelength, the equivalent expression is
\begin{equation}
 \Gamma_{X,q}(s)=4\pi\int_0^{\lambda_{\mathrm{th},q}}
 J_\lambda(s)\frac{\lambda}{hc}\,\sigma_{X,q}(\lambda)\,\mathrm{d}\lambda,
 \qquad \lambda_{\mathrm{th},q}=\frac{hc}{\chi_{X,q}}.
 \label{eq:ion-photo-wavelength}
\end{equation}
The spectral density and integration variable must use consistent units.
For example, $J_\lambda$ expressed per nanometre is integrated over
wavelength in nanometres, with the photon energy converted consistently.
Cross-sections vanish on the long-wavelength side of the threshold.
The integration grid must resolve the threshold and the spectral features
that contribute appreciably to the rate.

The analytic fits of \citet{Verner1996} provide a starting point for
ground-state cross-sections within their stated coverage. In a multilevel
model, excited levels have their own populations, thresholds, and
cross-sections. A ground-state rate then requires a population assumption
before it can represent the entire charge state.

\subsection{The radiation field and the role of opacity}

Radiation can reach a chromospheric parcel from the photosphere, the
transition region and corona, or emission within the chromosphere itself.
Photons from hotter regions can ionise atoms in a cooler parcel whose
local thermal electrons rarely exceed the threshold. The dominant source
depends on the element, wavelength, and position.

The specific intensity $I_\nu(s,\Omega)$ describes radiation arriving from a
direction $\Omega$. Its angular average is
\begin{equation}
 J_\nu(s)=\frac{1}{4\pi}\int I_\nu(s,\Omega)\,\mathrm{d}\Omega.
 \label{eq:ion-mean-intensity}
\end{equation}
Irradiance and mean intensity are different angular moments of the radiation
field. Converting an irradiance spectrum into $J_\nu$ therefore requires an
assumption about the emitting geometry.

For a uniformly bright photospheric surface, immediately above the surface,
the upward hemisphere gives $J_\nu=I_{\nu,\mathrm{ph}}/2$ and the emergent
surface flux is $F_{\nu,\mathrm{surf}}=\pi I_{\nu,\mathrm{ph}}$. Under this
assumption, an irradiance measured at distance $d$ gives
\begin{equation}
 J_{\nu,\mathrm{ph},0}
 =\frac{F_{\nu,d}}{2\pi}\left(\frac{d}{R_\odot}\right)^2.
 \label{eq:ion-irradiance-conversion}
\end{equation}
Downward coronal illumination requires its own geometry and emission model.
A spectrum measured at Earth does not uniquely determine the illumination
incident on a particular loop.

Along a ray with path coordinate $\ell$, radiative transfer satisfies
\begin{equation}
 \frac{\mathrm{d}I_\nu}{\mathrm{d}\ell}=-\kappa_\nu I_\nu+j_\nu,
 \label{eq:ion-transfer}
\end{equation}
where $\kappa_\nu$ is the extinction coefficient in $\mathrm{cm^{-1}}$ and
$j_\nu$ is the emission coefficient. Scattering contributes to extinction
and to the radiation redistributed into other directions and frequencies.
In the pure-absorption limit with no emission along the ray,
\begin{equation}
 I_\nu(\ell)=I_{\nu,0}\exp[-\tau_\nu(\ell)],
 \qquad
 \tau_\nu(\ell)=\int_{\ell_0}^{\ell}\kappa_\nu(\ell')\,\mathrm{d}\ell'.
 \label{eq:ion-attenuation}
\end{equation}

For ground-state neutral-hydrogen continuum absorption alone,
\begin{equation}
 \kappa_\nu=n_{\mathrm{H^0}}\sigma_{\mathrm{H^0}}(\nu),
 \qquad \tau_\nu=\sigma_{\mathrm{H^0}}(\nu)
          \int n_{\mathrm{H^0}}\,\mathrm{d}\ell.
 \label{eq:ion-h-opacity}
\end{equation}
Here the use of the total neutral density assumes that the ground state
dominates the absorbing neutral-H population.
Opacity determines the attenuation of an existing radiation field.
Calculating $\Gamma_X$ also requires the incident radiation or internal
emission and the target atom's cross-section.

Ground-state H continuum absorption vanishes below its ionisation threshold
of approximately $13.6\,\mathrm{eV}$. Metal continua and line opacity can still
affect radiation at these lower photon energies. In particular, resonant
absorption and scattering of Ly$\alpha$ require a line-transfer treatment
or a separately specified approximation. They are absent from
Equation~\eqref{eq:ion-h-opacity}. The assumed radiation and opacity treatment
should therefore accompany every calculated photoionisation rate
\citep{Laming2015,Dufresne2025}.

\subsection{Charge exchange with hydrogen}

Charge exchange transfers a bound electron between colliding particles:
\begin{equation}
 X^0+\mathrm{H^+}\rightleftharpoons X^++\mathrm{H^0}.
 \label{eq:ion-cx}
\end{equation}
In the forward direction, hydrogen captures an electron and $X$ becomes
ionised. In the reverse direction, $X^+$ captures hydrogen's electron and
becomes neutral. The electron remains bound to one of the reaction products,
and the total electric charge is conserved.

The corresponding rates per target particle are
\begin{equation}
 I_{\mathrm{cx},X}=n_{\mathrm{H^+}}k_{\mathrm{cx,ion},X},
 \qquad R_{\mathrm{cx},X}=n_{\mathrm{H^0}}k_{\mathrm{cx,rec},X}.
 \label{eq:ion-cx-rates}
\end{equation}
The coefficients have units of $\mathrm{cm^3\,s^{-1}}$ and depend on the
reaction and collision conditions. The appropriate heavy-particle
temperature need not equal $T_e$ in a plasma with different species
temperatures.

Oxygen and hydrogen have closely spaced first-ionisation thresholds.
Their nearly resonant charge exchange can consequently couple their
ionisation fractions when it dominates the other channels.
The fractions need not be exactly equal. Charge transfer can also affect
other elements, including sulphur, and requires species-specific data
\citep{Dufresne2021,Laming2021}. Elastic momentum-transfer collisions
enter the fractionation calculation separately; their collision frequencies
do not replace the reaction rates in Equation~\eqref{eq:ion-cx-rates}.

\subsection{Statistical equilibrium between neutral and singly ionised states}

For a two-stage model, the ionisation and neutralisation rates per particle are
\begin{align}
 I_X &= \Gamma_X+n_e C_X+n_{\mathrm{H^+}}k_{\mathrm{cx,ion},X},
 \label{eq:ion-forward-total}\\
 R_X &= n_e\alpha_X+n_{\mathrm{H^0}}k_{\mathrm{cx,rec},X}.
 \label{eq:ion-reverse-total}
\end{align}
Here $C_X$ describes $X^0\rightarrow X^+$, while $\alpha_X$ describes
$X^+\rightarrow X^0$. Statistical equilibrium requires
\begin{equation}
 n_{X^0}I_X=n_{X^+}R_X,
\end{equation}
which gives
\begin{equation}
 \xi_X=\frac{I_X}{I_X+R_X}.
 \label{eq:ion-two-state-equilibrium}
\end{equation}
This expression includes photoionisation, electron collisions, and charge
exchange. The term \emph{photoionisation equilibrium} is often used for
calculations including radiation, although it does not by itself establish
which channel dominates. The separate rates in
Equations~\eqref{eq:ion-forward-total} and \eqref{eq:ion-reverse-total}
allow that comparison to be made explicitly.

In the collision-only, two-stage limit,
\begin{equation}
 \xi_X^{\mathrm{CIE}}=\frac{C_X}{C_X+\alpha_X}.
 \label{eq:ion-two-state-cie}
\end{equation}
The common factor $n_e$ cancels when both coefficients depend only on
temperature. A multistage coronal calculation instead solves the coupled
population equations and returns the neutral fraction $f_{X,0}$.
Its total ionised fraction remains $1-f_{X,0}$.

\subsection{The LTE limit}

Under local thermodynamic equilibrium (LTE), the neutral-to-singly-ionised
ratio follows the Saha equation,
\begin{equation}
 \frac{n_{X^+}}{n_{X^0}}
 =\frac{2}{n_e}
  \left(\frac{2\pi m_e k_{\mathrm B}T}{h^2}\right)^{3/2}
  \frac{U_{X^+}(T)}{U_{X^0}(T)}
  \exp\left(-\frac{\chi_X}{k_{\mathrm B}T}\right),
 \label{eq:ion-saha}
\end{equation}
where $U(T)$ denotes a partition function. Replacing the partition functions
with constant statistical weights is an additional approximation.
Its adequacy depends on the population of excited levels.

LTE is a thermodynamic assumption, while statistical equilibrium specifies
that the population rates balance. Statistical equilibrium can hold in
non-LTE conditions when the radiation field and populations depart from their
local thermodynamic values. Recovering LTE from a rate solver requires the
relevant inverse processes and consistent detailed balance. An optically thin
ionisation and spontaneous-recombination network alone does not establish
the Saha limit.

\subsection{Extension to multiple stages and time dependence}

Let $I_q$ be the transition rate from stage $q$ to $q+1$, and $R_q$ the rate
from stage $q$ to $q-1$. With neighbouring-stage transitions, the local
population equations are
\begin{equation}
 \frac{\mathrm{d}f_{X,q}}{\mathrm{d}t}
 = I_{q-1}f_{X,q-1}+R_{q+1}f_{X,q+1}
   -(I_q+R_q)f_{X,q}.
 \label{eq:ion-multistage}
\end{equation}
Terms outside the retained stages are omitted, with $R_0=0$ and $I_Z=0$.
Statistical equilibrium sets the left-hand side to zero and enforces
Equation~\eqref{eq:ion-fractions}. A prescribed-background calculation holds
the atmospheric $T_e$, $n_e$, hydrogen populations, and specified radiation
fixed while solving the trace-element populations. Solving for $n_e$ through
charge conservation is a further coupling.

If conditions change faster than the populations relax, the time-dependent
equations must be evolved. For constant two-stage rates, the relaxation time
is $(I_X+R_X)^{-1}$. More general networks have several relaxation modes.
Hydrogen can retain a history of heating and shocks, which also changes
the collision partners and opacity seen by the metals \citep{Carlsson2002}.
Time-dependent calculations require initial populations and a sufficiently
resolved atmospheric and radiation history. Advection introduces transport
and boundary conditions in addition to the local reactions.

\subsection{Atomic coefficients from CHIANTI through fiasco}

The \texttt{fiasco} interface provides access to CHIANTI electron-impact
ionisation and recombination coefficients \citep{FiascoDocumentation}.
For a neutral--singly-ionised calculation, the required pair is the
\texttt{ionization\_rate} of the neutral ion object and the
\texttt{recombination\_rate} of the singly ionised object.
For example, Fe~I ionisation describes Fe~I to Fe~II, while Fe~II recombination
describes Fe~II to Fe~I. Fe~II ionisation instead describes Fe~II to Fe~III.
The two properties of a single ion object therefore refer to different
neighbouring-stage pairs.

The properties \texttt{direct\_ionization\_rate} and
\texttt{excitation\_autoionization\_rate} expose the ionisation components.
The properties \texttt{radiative\_recombination\_rate} and
\texttt{dielectronic\_recombination\_rate} expose recombination components.
For production calculations, the preferred total
\texttt{recombination\_rate} should be used because some ions have dedicated
total fits in CHIANTI \texttt{trparams} files.

Dataset availability must be checked for each ion. In the inspected
CHIANTI~11.0.2 Fe~II dataset, the \texttt{easplups} file is absent.
With fiasco~0.8.2, requesting the separate excitation--autoionisation
coefficient raises \texttt{MissingDatasetException}. Its total-ionisation
property handles this missing component internally and returns the
available direct-ionisation contribution. This is the database/software
treatment of the missing component; it does not demonstrate that the
physical process is identically zero.

Missing components should be recorded explicitly when plotting or exporting
the rates. The ordinary temperature-dependent properties also require
separate additions for the local photoionisation field, charge transfer,
and density effects appropriate to the selected model
\citep{Dufresne2024,Dufresne2025}.
